\subsection{Dasar Pemilihan Konfigurasi Uji}

Konfigurasi pengujian dalam penelitian ini disusun secara bertahap untuk memisahkan pengaruh penempatan otorisasi di \textit{edge}, penggunaan \textit{cache atribut}, strategi \textit{freshness}, dan mekanisme optimasi konfigurasi. \textit{Cloud-only} ABAC digunakan sebagai \textit{baseline} dengan verifikasi atribut terpusat karena sebagian besar sistem IAM pada \textit{edge} IoT masih bergantung pada verifikasi izin berbasis \textit{cloud}. \textcite{wang_edge-enabled_2025} menunjukkan bahwa banyak sistem \textit{edge} masih melakukan autentikasi dan otorisasi melalui IAM berbasis \textit{cloud} sehingga setiap komunikasi \textit{edge} tetap perlu diverifikasi ke \textit{cloud}. Pendekatan tersebut dapat meningkatkan ketergantungan terhadap koneksi \textit{remote} dan membatasi kemampuan sistem dalam menyediakan manajemen akses berlatensi rendah pada perangkat \textit{edge} IoT.

\textit{Edge-fresh} ABAC digunakan untuk mengevaluasi dampak pemindahan proses autorisasi ke \textit{edge} tanpa melibatkan \textit{cache} atribut sehingga pengaruh penempatan \textit{edge} dapat dipisahkan dari pengaruh \textit{caching}. \textcite{kalaria_adaptive_2024} menunjukkan bahwa kontrol akses berbasis \textit{fog} membawa pengambilan keputusan dan informasi kebijakan lebih dekat ke \textit{end nodes} untuk meningkatkan kemampuan adaptasi, ketersediaan sumber daya, dan mengurangi \textit{overhead} waktu pemrosesan pada lingkungan IoT yang dinamis. Pendekatan ini juga dirancang untuk mengurangi masalah \textit{overhead} komunikasi dan \textit{single point of failure} yang umum muncul pada arsitektur kontrol akses terpusat berbasis \textit{cloud}. Dengan demikian, konfigurasi \textit{edge-fresh} diperlukan untuk membedakan manfaat yang berasalh dari otorisasi \textit{edge} itu sendiri dengan manfaat tambahan yang berasal dari \textit{cache} atribut.

\textit{Static cache} ABAC digunakan sebagai \textit{baseline caching} karena \textit{context caching} merupakan pendekatan yang umum digunakan untuk meningkatkan efisiensi respons sistem IoT yang harus menangani \textit{context queries} secara \textit{real-time}. \textcite{weerasinghe_adaptive_2023} menunjukkan bahwa \textit{adaptive context caching} dapat meningkatkan efisiensi performa dan efisiensi biaya pada platform manajemen konteks ketika merespons \textit{context queries} dalam waktu mendekati \textit{real-time}. Namun, \textcite{weerasinghe_adaptive_2023} juga menekankan bahwa \textit{context caching} merupakan suatu siklus proses yang melibatkan \textit{selection, refreshing, eviction,} dan \textit{refresh overhead} yang dapat meningkatkan biaya pengelolaan konteks secara signifikan apabila tidak dikontrol dengan baik. \textcite{medvedev_refresh_2023} juga menjelaskan bahwa \textit{freshness, accuracy, completeness,} dan latensi merupakan QoS yang perlu dipertimbangkan bersama biaya operasional dalam manajemen konteks IoT. Oleh karena itu, \textit{static cache} ABAC digunakan untuk mengevaluasi \textit{trade-off} dasar antara pengurangan latensi dan risiko \textit{stale attribute} sebelum mekanisme \textit{freshness-aware} diterapkan.

\textit{Adaptive} TTL ABAC digunakan untuk mengevaluasi apakah mekanisme \textit{freshness-aware caching} dapat menghasilkan \textit{trade-off} yang lebih baik dibandingkan \textit{static caching}. \textcite{weerasinghe_adaptive_2023} menunjukkan bahwa \textit{context refreshing} yang terlalu sering dapat menyebabkan adanya lonjakan biaya manajemen konteks akibat meningkatnya \textit{retrieval} dan \textit{overhead} pemrosesan, sedangkan \textit{refresh} yang terlalu jarang dapat menurunkan validitas konteks yang tersimpan di \textit{cache}. \textcite{medvedev_refresh_2023} juga menunjukkan bahwa konfigurasi kecepatan \textit{refresh} memengaruhi keseimbangan antara biaya operasional dan QoS pada IoT \textit{middleware}. Dengan demikian, \textit{adaptive} TTL ABAC digunakan sebagai \textit{baseline} heuristik untuk menguji apakah TTL yang disesuaikan berdasarkan karakteristik atribut dapat mengurangi \textit{stale cache hit} tanpa menghasilkan \textit{overhead retrieval} yang berlebihan.

\textit{Random search} ABAC digunakan sebagai \textit{baseline} optimasi untuk mengontrol klaim bahwa peningkatan performa benar-benar berasal dari mekanisme optimasi, bukan sekadar akibat eksplorasi ruang konfigurasi. \textcite{bergstra_random_2012} menunjukkan bahwa \textit{random search} dapat menjadi \textit{baseline} yang kuat dalam optimisasi \textit{hyper-parameter} karena pada banyak kasus hanya sebagian kecil parameter yang memiliki pengaruh dominan terhadap performa sistem. Selain itu, \textcite{feng_dynamic_2025} menggunakan \textit{random search} sebagai \textit{baseline} pembanding terhadap GWO dan GA pada skenario optimisasi multi objektif. Oleh karena itu, \textit{random search} ABAC digunakan untuk mengevaluasi apakah optimasi metaheuristik benar-benar mampu menemukan konfigurasi TTL, ukuran \textit{cache}, dan \textit{cacheable attribute} yang lebih efektif dibandingkan pemilihan konfigurasi secara acak.

\textit{GWO-optimized} ABAC digunakan sebagai metode optimasi utama karena konfigurasi \textit{cache} atribut pada \textit{edge} ABAC membentuk masalah optimasi multi-parameter dengan \textit{trade-off} yang saling bertentangan antara latensi, \textit{freshness, overhead} komunikasi dan kualitas keputusan otorisasi. \textcite{mirjalili_grey_2014} menjelaskan bahwa GWO menggunakan mekanisme pencarian berbasis hirarki untuk mengeksplorasi ruang pencarian secara iteratif tanpa membutuhkan model analitik tertutup. \textcite{feng_dynamic_2025} menunjukkan bahwa GWO dapat diterapkan pada skenario optimisasi multi-objektif dan dibandingkan terhadap \textit{random search} serta GA pada ruang pencarian yang kompleks. Selain itu, \textcite{alsadie_modified_2025} menunjukkan bahwa modifikasi GWO pada lingkungan \textit{fog-cloud} mampu meningkatkan kemampuan eksplorasi dan eksploitasi sehingga menghasilkan penjadwalan dengan \textit{makespan} dan konsumsi energi yang lebih baik. Meskipun penelitian sebelumnya lebih banyak berfokus pada \textit{task schedulling} dan alokasi sumber daya, pendekatan tersebut menunjukkan bahwa GWO sesuai digunakan pada permasalahan optimasi \textit{edge/fog} yang memiliki ruang pencarian besar dan \textit{trade-off} multi-parameter. Oleh karena itu, \textit{GWO-optimized} ABAC digunakan untuk mengevaluasi apakah optimasi metaheuristik dapat meningkatkan keseimbangan antara performa \textit{cache} dan kualitas keputusan otorisasi dibandingkan pendekatan \textit{baseline} lainnya.

Dengan konfigurasi pengujian yang bertahap, hasil eksperimen tidak hanya digunakan untuk menunjukkan metode dengan performa terbaik, tetapi juga untuk mengisolasi kontribusi masing-masing komponen terhadap perubahan latensi, \textit{PIP calls, stale cache hit, false allow, false deny,} dan \textit{revocation false allow}. Pendekatan ini memungkinkan evaluasi dilakukan secara lebih beralasan dan terukur sehingga \textit{trade-off} antara performa, \textit{freshness}, dan reliabilitas keputusan otorisasi dapat dianalisis secara lebih jelas.